%% =====================================================================
%%  Bibliography.
%%
%%  DISCIPLINE.  Every entry below is either (a) resolved against a source
%%  fetched and read during the work reported here (the arXiv/journal record,
%%  or the text of the cited paper itself), or (b) marked UNVERIFIED-... and
%%  listed as such.  Nothing was invented.  Where only part of an entry could
%%  be confirmed, the entry carries a bracketed note saying which part.
%%
%%  REFEREE pass: the single UNVERIFIED entry (Lai-Lupu-Sprang, arXiv:2505.23088)
%%  has been resolved -- title and author initials from the arXiv API and the
%%  arXiv abstract page, journal reference from the Crossref record for
%%  DOI 10.1007/s40687-025-00559-x -- and the key renamed UNVERIFIED-LLS2025 ->
%%  LLS2025 throughout.  Bel2019, RhinViola, Beukers2008, KubotaLeopoldt1964 and
%%  FSZ2019 were re-checked against fetched records and are exact as printed; in
%%  particular FSZ2019 does NOT carry the Compos. Math. 154 (2018) error found in
%%  two sibling papers.  No entry now carries an UNVERIFIED marker.
%% =====================================================================
\begin{thebibliography}{99}

\bibitem{Apery1979}
R.~Ap\'ery,
\emph{Irrationalit\'e de $\zeta(2)$ et $\zeta(3)$},
in: Journ\'ees Arithm\'etiques (Luminy, 1978),
Ast\'erisque \textbf{61} (1979), 11--13.

\bibitem{BallRivoal2001}
K.~Ball and T.~Rivoal,
\emph{Irrationalit\'e d'une infinit\'e de valeurs de la fonction z\^eta aux
entiers impairs},
Invent.\ Math.\ \textbf{146} (2001), no.~1, 193--207.

\bibitem{Bel2019}
P.~Bel,
\emph{Irrationalit\'e des valeurs de $\zeta_p(4,x)$},
J.\ Th\'eor.\ Nombres Bordeaux \textbf{31} (2019), no.~1, 81--99.
(Lemme~3.2 is the $p$-adic irrationality criterion used throughout; it is quoted
as Lemma~4 of \cite{LSZ2025}.)

\bibitem{Beukers2008}
F.~Beukers,
\emph{Irrationality of some $p$-adic $L$-values},
Acta Math.\ Sin.\ (Engl.\ Ser.)\ \textbf{24} (2008), no.~4, 663--686;
arXiv:math/0603277.
(Theorem~11.2 carries the $\zeta_2(3)$ and $\zeta_3(3)$ ledgers;
Proposition~11.1(3) is the coefficient-growth bound; Proposition~5.2 gives
$T_2(1/4)=4^3\zeta_2(3)$, $T_3(1/3)=3^3\zeta_3(3)$ and the $p=5$ pair;
Corollary~11.3 is the irrationality of the individual Hurwitz values.)

\bibitem{Brown2014}
F.~Brown,
\emph{Irrationality proofs for zeta values, moduli spaces and dinner parties},
arXiv:1412.6508.

\bibitem{Brown2026}
F.~Brown,
\emph{Mellin transforms, transfinite diameter and rational approximations of
integrals},
arXiv:2604.20741.
(Remark~47 records the referee's question about Rhin--Viola prime cancellation.)

\bibitem{BrownZudilin2022}
F.~Brown and W.~Zudilin,
\emph{On cellular rational approximations to $\zeta(5)$},
arXiv:2210.03391.

\bibitem{Calegari2005}
F.~Calegari,
\emph{Irrationality of certain $p$-adic periods for small $p$},
Int.\ Math.\ Res.\ Not.\ (2005), no.~20, 1235--1249;
arXiv:math/0408214.

\bibitem{FSZ2019}
S.~Fischler, J.~Sprang and W.~Zudilin,
\emph{Many odd zeta values are irrational},
Compositio Math.\ \textbf{155} (2019), no.~5, 938--952; arXiv:1803.08905.

\bibitem{KubotaLeopoldt1964}
T.~Kubota and H.~W.~Leopoldt,
\emph{Eine $p$-adische Theorie der Zetawerte. Teil~I: Einf\"uhrung der
$p$-adischen Dirichletschen $L$-Funktionen},
J.\ Reine Angew.\ Math.\ \textbf{214/215} (1964), 328--339.

\bibitem{Lai2024}
L.~Lai,
\emph{Small improvements on the Ball--Rivoal theorem and its $p$-adic variant},
arXiv:2407.14236.
(Sections 7--8 contain the multi-parameter Volkenborn family $V_n(t)$ and the
normalising constant $p^{(l+1/(p-1))M\varphi(p^l)n}$.)

\bibitem{Lai2025}
L.~Lai,
\emph{On the irrationality of certain $2$-adic zeta values},
Int.\ J. Number Theory \textbf{21} (2025), no.~1, 207--235;
arXiv:2304.00816.
(The families $A_n$ and $B_n$ are (3.1) and (3.2); the $\triangle$-operator
estimate is Lemma~2.4, quoted as Lemma~6 of \cite{LSZ2025}.)

\bibitem{LaiSprang2023}
L.~Lai and J.~Sprang,
\emph{Many $p$-adic odd zeta values are irrational},
arXiv:2306.10393.
(Lemma~12 is the coset decomposition \eqref{eq:cosetsum}; Lemma~21 is the
Volkenborn linear-form identity used in Theorem~\ref{thm:transfer}(ii);
Definition~16 is the general rational function.)

\bibitem{LSZ2025}
L.~Lai, J.~Sprang and W.~Zudilin,
\emph{A note on the irrationality of $\zeta_2(5)$},
arXiv:2505.05005.
(Definitions 10--11; Lemma~4 = Bel's criterion; Lemma~5 = the Volkenborn
translation formula; Lemmas 6--7 = the $\triangle$-operator; Lemma~9 = the
$\theta=\tfrac12$ identity; Lemma~15(b) = the Casoratian; Lemma~20 = the
coefficient-growth bound; Remark~13 and the Final remarks are quoted in
\S\ref{ss:questions} and \S\ref{ss:principle}.)

\bibitem{LLS2025}
L.~Lai, C.~Lupu and J.~Sprang,
\emph{On the irrationality of certain $p$-adic zeta values},
Res.\ Math.\ Sci.\ \textbf{12} (2025), no.~4, Paper No.~77;
arXiv:2505.23088.
(Definition~3.1 is the $p\ge5$ construction
$R_n(t)=p^{pn}n!^{s}t^{M_0}\prod_{j=1}^{p-1}(t+j/p)_n/(t)_{n+1}^{p-1+s}$;
Theorem~1.1 is the constant $c_p$ of Finding~\ref{find:LLS} and
Corollary~\ref{cor:noAhelp}.)

\bibitem{RhinViola}
G.~Rhin and C.~Viola,
\emph{On a permutation group related to $\zeta(2)$},
Acta Arith.\ \textbf{77} (1996), no.~1, 23--56;
\emph{The group structure for $\zeta(3)$},
Acta Arith.\ \textbf{97} (2001), no.~3, 269--293.

\bibitem{Zudilin2004}
W.~Zudilin,
\emph{Arithmetic of linear forms involving odd zeta values},
J.\ Th\'eor.\ Nombres Bordeaux \textbf{16} (2004), no.~1, 251--291;
arXiv:math/0206176.
(Lemma~8 is the realisation of the group used in \S\ref{ss:groupexp};
Lemmas~7--9 and equation (4.4) give the invariant \eqref{eq:invariant};
Lemma~12 is the saddle-point constant $C_0=-f_0(\tau_0)$; Lemma~16 is the
denominator lemma; equation (2.4) is the unequal-length form.)

\end{thebibliography}
